\documentclass[11pt,a4paper]{article}

% ── Packages ──────────────────────────────────────────────────────────────────
\usepackage[margin=2.5cm]{geometry}
\usepackage{amsmath,amssymb}
\usepackage{graphicx}
\usepackage{booktabs}
\usepackage{longtable}
\usepackage{array}
\usepackage{xcolor}
\usepackage{hyperref}
\usepackage{caption}
\usepackage{subcaption}
\usepackage{enumitem}
\usepackage{microtype}
\usepackage{listings}
\usepackage{float}
\usepackage{placeins}

\graphicspath{{figures/}}

\hypersetup{
  colorlinks = true,
  linkcolor  = blue!70!black,
  citecolor  = green!50!black,
  urlcolor   = blue!70!black,
}

% ── Custom commands ───────────────────────────────────────────────────────────
\newcommand{\numu}{$\nu_\mu$}
\newcommand{\numubar}{$\bar{\nu}_\mu$}
\newcommand{\uboone}{MicroBooNE}
\newcommand{\GeVc}{\,\text{GeV}\!/c}
\newcommand{\pot}{\,\text{POT}}
\newcommand{\cm}{\,\text{cm}}
\newcommand{\Npi}{$N\pi^\pm$}

\lstset{
  language=C++,
  basicstyle=\ttfamily\small,
  keywordstyle=\bfseries\color{blue!70!black},
  commentstyle=\itshape\color{gray},
  backgroundcolor=\color{gray!10},
  frame=single,
  framerule=0.4pt,
  numbers=none,
  breaklines=true,
}

\newcolumntype{C}[1]{>{\centering\arraybackslash}p{#1}}
\newcolumntype{R}[1]{>{\raggedleft\arraybackslash}p{#1}}
\newcolumntype{L}[1]{>{\raggedright\arraybackslash}p{#1}}
\emergencystretch=3em

% ─────────────────────────────────────────────────────────────────────────────
\begin{document}

\begin{titlepage}
  \centering
  \vspace*{2cm}
  {\LARGE\bfseries Measurement of $\nu_\mu$ Charged-Current Multi-Pion Production
  on Argon with the MicroBooNE Detector in the NuMI Beam\par}
  \vspace{0.8cm}
  {\large Technical Note --- Two- and Three-Charged-Pion Final States
  (\texttt{CC1mu2pi}, \texttt{CC1mu3pi})\par}
  \vspace{0.5cm}
  {\normalsize Draft 0.1 --- 2026-08-13\par}
  \vspace{1.5cm}
  % NOTE: author list to be confirmed. Contributors to the multi-pion effort:
  {\large Tusharadri Mahmud, Linhui Gu, Ishanne Pophale, Jarek Nowak\par}
  \vspace{0.5cm}
  {\normalsize\textit{Lancaster University}\par}
  \vspace{2cm}
  \begin{abstract}
    We describe the extension of the \uboone{} NuMI $\nu_\mu$ charged-current
    single-charged-pion analysis to higher pion multiplicities: charged-current
    interactions producing a muon and exactly two or three charged pions (and any
    number of protons) in the final state. These \Npi{} topologies probe the
    resonant and deep-inelastic regimes and the intranuclear pion cascade
    (final-state interactions), and are a significant background to the single-pion
    signal. The selections are built as thin subselections of the inclusive
    single-pion selection, adopting the looser pion identification of the dedicated
    \uboone{} multi-pion study to recover efficiency. On the Run-1 forward-horn-current
    $\nu_\mu$ overlay the two-pion (three-pion) selection reaches a signal efficiency
    of $15.5\%$ ($6.1\%$) at an intrinsic purity of $23\%$ ($12\%$), corresponding to
    of order $580$ ($40$) selected signal events at full exposure. We present the
    selection, the cut-flow, and a study of the pion-momentum thresholds that set the
    accessible signal phase space; the two-pion channel supports a flux-averaged total
    cross section and a coarse differential measurement, while the three-pion channel
    is a total-cross-section measurement. This note is a companion to the single-pion
    analysis note, from which the detector model, samples, exposure, systematic
    framework and unfolding are inherited unchanged.
  \end{abstract}
\end{titlepage}

\tableofcontents
\newpage

% =============================================================================
\section{Introduction}
\label{sec:intro}

Charged-current (CC) production of multiple charged pions on nuclei sits at the
transition between the resonant (RES) and deep-inelastic (DIS) regimes of
neutrino--nucleus scattering and is strongly shaped by the intranuclear hadronic
cascade. A $\nu_\mu$ interaction that produces several pions at the primary vertex may
have some of them absorbed, charge-exchanged or produced in re-scattering before they
leave the nucleus, so the observed charged-pion multiplicity is a direct probe of
final-state interactions (FSI). Multi-pion final states are also a non-negligible
background to exclusive single-pion and pion-less measurements, so a dedicated
measurement of their rate and kinematics on argon is valuable both as a cross-section
result and as a constraint on the interaction model used across the \uboone{} programme.

This note documents the extension of the NuMI $\nu_\mu\,\mathrm{CC}\,1\pi^\pm$ analysis
(the ``single-pion note'', to which this is a companion) to the two- and
three-charged-pion final states. The selections, \texttt{CC1mu2pi} and
\texttt{CC1mu3pi}, are implemented as subselections of the inclusive single-pion
selection \texttt{CC1mu1piXp} in the \texttt{xsec\_analyzer} framework and reuse its
muon reconstruction, samples, per-run POT normalisation, systematic universes and
Wiener-SVD unfolding without modification. The multi-pion pion identification follows
the dedicated \uboone{} multi-pion study of
Mahmud~\cite{mahmud} and the accompanying analyzer of Gu~\cite{gu}.

Given the steeply falling multi-pion cross section and the reduced reconstruction
efficiency for events with several (often soft) pions, these channels are
statistically limited: the emphasis of this first note is the selection definition,
its efficiency and purity, the accessible phase space, and a feasibility assessment
for total and differential cross sections.

% =============================================================================
\section{Signal definition}
\label{sec:signal}

The signal is a $\nu_\mu$ (or $\bar\nu_\mu$) charged-current interaction with a
true muon and \emph{exactly} $N$ charged pions ($\pi^+$ or $\pi^-$; \uboone{} does not
determine the pion charge sign) in the final state, with $N=2$ (\texttt{CC1mu2pi}) or
$N=3$ (\texttt{CC1mu3pi}). The signal is inclusive with respect to protons, neutrons and
neutral pions/electromagnetic activity, subject to the requirement that there is exactly
one muon and no other charged leptons. The interaction vertex must lie in the fiducial
volume defined in the single-pion note.

\paragraph{Pion-momentum threshold.} The accessible signal phase space is bounded from
below by the momenta of the pions themselves: a pion must be energetic enough to produce
a reconstructable track. In a genuine two- or three-pion event the sub-leading pions are
systematically softer than the leading one, and the softest pion sets the floor on how
often the full topology can be reconstructed. The single-pion signal definition applies
a threshold to \emph{the} pion; for the multi-pion channels this must be generalised so
that the momentum requirement applies consistently to \emph{all} $N$ pions rather than to
one. Section~\ref{sec:thresholds} quantifies the true leading and sub-leading pion
spectra and studies the selection efficiency as a function of a per-pion momentum
threshold, and recommends the threshold that defines the measurable phase space. That study
(Section~\ref{sec:thresholds}) adopts a per-pion threshold of $p_\pi>0.10\GeVc$ applied
to \emph{all} $N$ pions (the pion tracking turn-on), together with the muon threshold
$p_\mu>0.15\GeVc$ inherited from the single-pion definition.

% =============================================================================
\section{Samples, exposure and normalisation}
\label{sec:samples}

The samples, beam exposure, per-run protons-on-target normalisation, beam-off (EXT)
cosmic treatment and dirt overlay are identical to the single-pion note and are not
repeated here. All results in this draft use the Run-1 forward-horn-current (FHC)
$\nu_\mu$ overlay for the efficiency and cut-flow studies; the extension to the full
FHC, RHC and combined exposures follows the same per-run machinery. As in the parent
analysis, the beam-on data remain blinded and all data-like distributions use
Poisson-thrown fake data.

% =============================================================================
\section{Event selection}
\label{sec:selection}

\texttt{CC1mu2pi} and \texttt{CC1mu3pi} inherit the complete \texttt{CC1mu1piXp}
selection --- the same slice, fiducial-volume, topology and muon-candidate
requirements --- and the same reconstructed observables. Only the pion identification
and the required charged-pion multiplicity differ. In the code the multiplicity is a
virtual (\texttt{required\_charged\_pions()}) so that the derived value ($N=2,3$) takes
effect inside the inherited signal definition and reco selection.

Re-using the single-pion pion identification directly gives a prohibitively low
efficiency ($4.2\%$ for $2\pi$, $0.8\%$ for $3\pi$), because that identification was
tuned to isolate a single energetic pion and reject protons: a tight $4\cm$ vertex
window, a track-multivariate/Bragg/length requirement, and cuts on the single
muon--pion opening angle and on electromagnetic showers. Following
Mahmud~\cite{mahmud} and Gu~\cite{gu}, the multi-pion selections loosen the pion
identification through five parameters, implemented as overridable virtuals on the base
class so that the single-pion selection is left byte-for-byte unchanged
(Table~\ref{tab:mp_params}):

\begin{table}[H]
  \centering
  \caption{Pion-identification parameters overridden by the multi-pion selections.
    The inclusive defaults reproduce the single-pion selection exactly.}
  \label{tab:mp_params}
  \small
  \begin{tabular}{lcc}
    \toprule
    Parameter & Single-pion default & Multi-pion \\
    \midrule
    Vertex-to-track window       & $4\cm$ & $9.5\cm$ \\
    Pion identification          & LLR-PID $+$ TMVA $+$ Bragg $+$ length & LLR-PID only \\
    Uncontained pions allowed    & $0$ & $1$ \\
    $\mu$--$\pi$ opening-angle cut & applied & not applied \\
    Shower (\,$\pi^0$) veto      & applied & not applied \\
    \bottomrule
  \end{tabular}
\end{table}

The wider vertex window recovers pions from secondary scatters and slightly displaced
starts; the log-likelihood-ratio-only identification drops the proton-rejection cuts that
were specific to the one-pion topology; the ``maximum uncontainment'' tolerance allows a
single pion per event to leave the detector; and the opening-angle and shower cuts, which
are one-pion-topology cuts (a single $\mu$--$\pi$ pair, and an electromagnetic veto), are
dropped because they are ill-defined or unwanted for a multi-pion, EM-inclusive signal.
In the pion loop the contained and uncontained pion candidates are tallied separately and
the charged-pion count is formed as the number of contained pions plus up to one
uncontained pion.

% =============================================================================
\section{Selection performance}
\label{sec:performance}

Table~\ref{tab:mp_cutflow} gives the cut-flow for both selections on the Run-1 FHC
$\nu_\mu$ overlay. The dominant efficiency loss is at the $N$-pion identification stage
(MuonCandidate$\rightarrow$ContainedPions): reconstructing $N$ pion-like tracks is the
efficiency ceiling, while the upstream vertex/topology/muon requirements and the
downstream wire-gap and track-multiplicity cuts are comparatively gentle. The loosened
identification raises the efficiency by a factor of $\sim\!4$ ($2\pi$) and $\sim\!8$
($3\pi$) relative to the single-pion identification, turning $3\pi$ from a
non-measurement ($\mathcal{O}(1)$ event) into a viable one ($\sim\!40$ events).

\begin{table}[H]
  \centering
  \caption{Multi-pion cut-flow yields (Run-1 FHC $\nu_\mu$ overlay, POT-weighted;
    beam-off cosmics and dirt not included). ``Total'' is signal $+$ MC background; Eff.\
    is the cumulative signal efficiency relative to the generated signal (None stage);
    Pur.\ $=$ Signal/Total and is MC-intrinsic (an upper bound --- the real-sample purity
    is lower once EXT cosmics, which dominate at these multiplicities, are included). The
    shower veto and opening-angle cut are not applied.}
  \label{tab:mp_cutflow}
  \small
  \begin{tabular}{lrrrrr}
    \toprule
    Cut & Total & Signal & Bkg MC & Eff. [\%] & Pur. [\%] \\
    \midrule
    \multicolumn{6}{l}{\textit{\texttt{CC1mu2pi} --- exactly 2 charged pions}} \\
    None            & 759\,496 & 3\,746 & 755\,750 & 100.0 & 0.5 \\
    InFiducialVol   & 177\,783 & 3\,145 & 174\,639 & 83.9 & 1.8 \\
    Topological     & 83\,662 & 2\,565 & 81\,097 & 68.5 & 3.1 \\
    MuonCandidate   & 58\,791 & 2\,273 & 56\,518 & 60.7 & 3.9 \\
    ContainedPions  & 3\,372 & 781 & 2\,591 & 20.9 & 23.2 \\
    MuonIn3Planes   & 3\,181 & 737 & 2\,444 & 19.7 & 23.2 \\
    PionIn3Planes   & 2\,786 & 652 & 2\,134 & 17.4 & 23.4 \\
    \textbf{Final}  & \textbf{2\,500} & \textbf{579} & \textbf{1\,921} & \textbf{15.5} & \textbf{23.2} \\
    \midrule
    \multicolumn{6}{l}{\textit{\texttt{CC1mu3pi} --- exactly 3 charged pions}} \\
    None            & 759\,496 & 656 & 758\,840 & 100.0 & 0.09 \\
    InFiducialVol   & 177\,783 & 558 & 177\,225 & 85.0 & 0.3 \\
    Topological     & 83\,662 & 457 & 83\,205 & 69.6 & 0.5 \\
    MuonCandidate   & 58\,791 & 393 & 58\,398 & 59.9 & 0.7 \\
    ContainedPions  & 447 & 53 & 394 & 8.1 & 11.8 \\
    MuonIn3Planes   & 415 & 50 & 365 & 7.6 & 12.0 \\
    PionIn3Planes   & 384 & 47 & 337 & 7.2 & 12.2 \\
    \textbf{Final}  & \textbf{344} & \textbf{40} & \textbf{304} & \textbf{6.1} & \textbf{11.6} \\
    \bottomrule
  \end{tabular}
\end{table}

These efficiencies sit below the values reported in the dedicated
study~\cite{mahmud} ($37\%$ for $2\pi$, $31\%$ for $3\pi$) because the selections here
retain the inherited muon-momentum-quality, wire-gap, track-multiplicity and
software-trigger requirements that the leaner study omits; the trade is a lower
efficiency for a markedly higher intrinsic purity, which is the more favourable operating
point for a cross-section extraction.

% =============================================================================
\section{Pion-momentum thresholds and the accessible phase space}
\label{sec:thresholds}

The efficiency ceiling at the $N$-pion identification stage is driven by the softest
pions in the event. The signal multiplicity is currently defined by counting charged
pions above a negligible momentum ($p_\pi>0$), so a true $N$-pion event is counted as
signal however soft its sub-leading pions are --- including pions too soft to produce a
reconstructable track. Such events enter the efficiency denominator but can never enter
the numerator, mechanically depressing the quoted efficiency. A physically meaningful
signal definition applies a momentum threshold to \emph{all} $N$ pions, restricting the
measurement to a phase space that can actually be reconstructed.

\paragraph{Method.} For each true $N$-pion signal event we record the leading (hardest)
and minimum (softest) true charged-pion momenta, \texttt{mc\_pionpm\_lead\_mom} and
\texttt{mc\_pionpm\_min\_mom}. The efficiency as a function of a per-pion threshold
$p_{\pi}^{\min}$ is then evaluated on the sub-sample whose \emph{softest} pion exceeds
the threshold,
\[
  \varepsilon(p_\pi^{\min}) =
  \frac{N\big(\text{selected}\ \wedge\ \text{signal}\ \wedge\ p_\pi^{\text{soft}}>p_\pi^{\min}\big)}
       {N\big(\text{signal}\ \wedge\ p_\pi^{\text{soft}}>p_\pi^{\min}\big)},
\]
which isolates the effect of removing unreconstructable soft-pion events from the signal
definition.

\paragraph{Results.} The sub-leading pions are indeed substantially softer than the
leading one: on the Run-1 FHC overlay the mean leading (softest) true charged-pion
momentum is $0.52$ ($0.27$)$\GeVc$ for $2\pi$ and $0.57$ ($0.21$)$\GeVc$ for $3\pi$
(Table~\ref{tab:mp_thr}). A significant fraction of events have a pion below the
tracking turn-on: $10\%$ ($18\%$) of $2\pi$ ($3\pi$) signal events have their softest
pion below $0.10\GeVc$, rising to $30\%$ ($51\%$) below $0.175\GeVc$. These events, with
an unreconstructable pion, are counted as signal at the current $p_\pi>0$ threshold but
can never pass the reco selection.

\begin{table}[H]
  \centering
  \caption{Signal efficiency as a function of a per-pion momentum threshold
    $p_\pi^{\min}$ (applied to all $N$ pions via the softest one), Run-1 FHC $\nu_\mu$
    overlay. ``Kept'' is the fraction of the $p_\pi>0$ signal retained; Eff.\ is the
    efficiency evaluated within the restricted phase space.}
  \label{tab:mp_thr}
  \small
  \begin{tabular}{lcccc}
    \toprule
    & \multicolumn{2}{c}{\texttt{CC1mu2pi}} & \multicolumn{2}{c}{\texttt{CC1mu3pi}} \\
    \cmidrule(lr){2-3}\cmidrule(lr){4-5}
    $p_\pi^{\min}$ [$\GeVc$] & Kept & Eff.\ [\%] & Kept & Eff.\ [\%] \\
    \midrule
    $0$ (pre-fix)  & $1.00$ & $16.3$ & $1.00$ & $6.4$ \\
    $0.10$        & $0.90$ & $17.1$ & $0.82$ & $6.2$ \\
    $0.175$       & $0.70$ & $18.9$ & $0.49$ & $8.1$ \\
    $0.25$        & $0.43$ & $20.2$ & $0.26$ & $5.6$ \\
    \bottomrule
  \end{tabular}
\end{table}

The central result is that raising the threshold \emph{barely changes the efficiency}:
restricting to events whose softest pion exceeds $0.175\GeVc$ --- i.e.\ where all pions
are comfortably above tracking turn-on --- still yields an efficiency of only $\sim\!19\%$
($2\pi$). Roughly four in five events with fully reconstructable pions are still lost. The
efficiency ceiling is therefore \emph{intrinsic multi-track reconstruction} --- resolving
$N$ distinct pion tracks from a common vertex in the presence of hadronic re-interactions
and track merging --- and not the soft-pion threshold. The soft sub-leading pions are a
real but secondary effect ($10$--$18\%$ of events genuinely below tracking threshold).

\paragraph{Recommendation.} A per-pion momentum threshold of $p_\pi^{\min}\simeq0.10\GeVc$
(the pion tracking turn-on) applied to all $N$ pions is adopted for the signal definition:
it removes the genuinely unreconstructable soft tail ($10$--$18\%$ of events) at negligible
efficiency cost, defines an honest measurable phase space, and preserves $\sim\!90\%$
($2\pi$) / $\sim\!82\%$ ($3\pi$) of the signal --- important given the limited statistics,
especially for $3\pi$. A tighter $0.175\GeVc$ threshold, matching the single-pion measured
phase space, would be cleaner but sacrifices half of the $3\pi$ signal and is not justified
by the marginal efficiency gain. The threshold does not, and is not expected to, recover
the bulk of the efficiency, which is limited by reconstruction and is corrected for by the
unfolding.

The threshold scan above was performed on the pre-fix signal definition, which also
contained a bookkeeping error: the momentum and opening-angle requirements were applied to
\texttt{TrueCandidatePionP}, the \emph{last} charged pion encountered in the daughter loop,
rather than to a well-defined pion. For the single-pion signal this is unambiguous (there
is only one pion), but for $N>1$ it applied the cuts to an arbitrary, ordering-dependent
pion. Both issues have now been fixed: the momentum threshold is applied to the softest of
the $N$ pions (so all $N$ are above threshold) and the opening angle to the leading pion.
The correction removed the spurious ordering-dependent $0.175\GeVc$ cut on an arbitrary
pion, so the well-defined signal count rises (e.g.\ $3\,251\rightarrow3\,746$ for $2\pi$)
and the efficiency settles at $15.5\%$ ($2\pi$) / $6.1\%$ ($3\pi$) in the corrected phase
space (Table~\ref{tab:mp_cutflow}), slightly below the scan projection because the fix
broadens the signal to include soft-leading and wide-angle events that the arbitrary cut
had randomly rejected. We verified that the single-pion selection is bit-for-bit unchanged
by the fix (identical signal, selected and efficiency on the Run-1 FHC overlay).

% =============================================================================
\section{Backgrounds}
\label{sec:backgrounds}

At these multiplicities the beam-off cosmic (EXT) contribution is expected to dominate the
background, followed by $\nu_\mu\,\mathrm{CC}$ events of the wrong charged-pion
multiplicity (single-pion events with an extra mis-identified track, or higher-multiplicity
events losing a pion) and neutral-current interactions with charged-pion-like tracks. The
cut-flow of Table~\ref{tab:mp_cutflow} is the $\nu_\mu$ overlay only and therefore does not
yet include EXT or dirt; a full background composition, including the data-driven EXT
subtraction inherited from the parent analysis, is deferred to the next revision. The MC
background composition and any dedicated multi-pion background-control sidebands will be
reported alongside.

% =============================================================================
\section{Feasibility and observables}
\label{sec:feasibility}

With of order $580$ selected signal events at full exposure, the two-pion channel can
support a flux-averaged total cross section and, plausibly, a coarse (two- to three-bin)
differential measurement in a robust observable such as the leading-pion momentum or the
total visible hadronic energy. The three-pion channel, with of order $40$ selected signal
events, is a flux-averaged total-cross-section measurement. Both will use the same
Wiener-SVD unfolding and systematic treatment as the parent analysis; the reduced
statistics argue for coarse binning and, for $3\pi$, a single integrated bin. A concrete
observable set and binning will be fixed once the threshold study defines the signal phase
space and the background model is in place.

% =============================================================================
\section{Systematic uncertainties}
\label{sec:systematics}

The systematic framework --- flux (PPFX), cross-section (GENIE reweight), reinteraction
and detector-variation universes, propagated through the response matrix and unfolding ---
is inherited unchanged from the single-pion note. The multi-pion channels are expected to
be statistics-limited, so a full systematic evaluation is deferred until the selection and
phase space are finalised; the detector-variation samples and the flux/cross-section
universes are already available in the same processed files.

% =============================================================================
\section{Status and follow-up}
\label{sec:status}

\begin{itemize}[nosep]
  \item Selections \texttt{CC1mu2pi} / \texttt{CC1mu3pi} implemented, compiled and
        validated on Run-1 FHC; cut-flow, efficiency and purity established
        (Section~\ref{sec:performance}).
  \item Pion-momentum threshold study in progress (Section~\ref{sec:thresholds}); will fix
        the accessible phase space and the per-pion threshold, and correct the signal-def
        bookkeeping for $N>1$.
  \item To do: full background composition including EXT and dirt; extension to full FHC,
        RHC and combined exposures; observable/binning choice; total (and, for $2\pi$,
        coarse differential) cross-section extraction with the inherited unfolding and
        systematics.
\end{itemize}

% =============================================================================
\begin{thebibliography}{9}
\bibitem{mahmud} T.\ Mahmud, \textit{Study of $\nu_\mu$ charged-current three-charged-pion
  production in \uboone{}} (\texttt{CC3}$\pi^+$ selection), MSc by Research thesis,
  Lancaster University.
\bibitem{gu} L.\ Gu, multi-pion event-selection analyzer for the \uboone{} NuMI
  cross-section framework (private communication / internal code).
\bibitem{note1} I.\ Pophale and J.\ Nowak, \textit{Measurement of
  $\nu_\mu+\bar\nu_\mu$ charged-current single-pion production differential cross sections
  on argon with \uboone{} in the NuMI beam}, \uboone{} analysis note (companion document).
\end{thebibliography}

\end{document}
